% !TeX root = ../../Book5.tex
% 纲要标题：### 31.3 余切复形的具体构造
% 纲要位置：几何学计划纲要.md，第 6042 行。
% 正文状态：未撰写；本文件目前是框架。

% 纲要内容开始
% 
% * 多项式呈示与两项截断
% 
%   从 A→P→B 构造 [I/I²→Ω_{P/A}⊗_P B]；它是余切复形的低阶截断模型，不在非完全交情形直接等同于完整 L_{B/A}。
% 
% * 局部完全交与幅度
% 
%   对由正则序列定义的商证明两项模型计算完整余切复形，并核验其幅度为 [−1,0]。
% 
% * 单纯多项式消解
% 
%   建立自由多项式代数的单纯消解、面与退化映射及所用比较同伦；先说明其与模的自由消解不同。
% 
% * 一般余切复形的构造
% 
%   对消解逐项取微分模并张量到 B，再以正规化复形形成 L_{B/A}；统一链与上链次数。
% 
% * 消解独立性与函子性
% 
%   通过共同消解和同伦比较证明导出对象独立于呈示，不仅核对 H⁰ 相同。
% 
% * 传递三角
% 
%   从 A→B→C 的相容消解构造三角，并验证微分模传递序列是其低阶后果。
% 
% * 导出基变换及普通基变换的条件
% 
%   把导出交换图与普通张量模型分开，注明 Tor 独立等可去掉导出的条件。
% 
% 纲要内容结束

\section{余切复形的具体构造}
\label{g5:sec:3103}

待撰写

\subsection{多项式呈示与两项截断}
\label{g5:subsec:3103:01}
待撰写

\subsection{局部完全交与幅度}
\label{g5:subsec:3103:02}
待撰写

\subsection{单纯多项式消解}
\label{g5:subsec:3103:03}
待撰写

\subsection{一般余切复形的构造}
\label{g5:subsec:3103:04}
待撰写

\subsection{消解独立性与函子性}
\label{g5:subsec:3103:05}
待撰写

\subsection{传递三角}
\label{g5:subsec:3103:06}
待撰写

\subsection{导出基变换及普通基变换的条件}
\label{g5:subsec:3103:07}
待撰写
% 纲要条目排为编号小节标题，并进入目录与书签；说明段仍保留于纲要注释中。
